% Generated by scripts/assemble.py; corrections belong in source/corrections.json.
\examhead{143}{湖南大学 2026 年数学分析试卷}{136}
% p137-b11; source PDF page 137
\blockanchor{p137-b11}
\question{1}试求极限 $\lim_{n\to \infty}\sum_{i = 1}^{n}\frac{\sin\frac{i}{n}\pi}{n + \frac{i}{n}}.$\par
% p137-b12; source PDF page 137
\blockanchor{p137-b12}
\question{2}设 $f(x)$ 在 $[0,1]$ 上连续, 在 $(0,1)$ 上二阶可导, $f(0)f(1) > 0$ 且 $f''(x) > 0$, $\int_0^1 f(x) \mathrm{d}x = 0$.\par
% p137-b13; source PDF page 137
\blockanchor{p137-b13}
(1) 证明: $f(x)$ 在 $(0,1)$ 上恰有两个零点.\par
% p137-b14; source PDF page 137
\blockanchor{p137-b14}
(2) 证明: 存在 $c \in (0,1)$, 使得 $f'(c) = \int_{0}^{c} f(x) \mathrm{d}x$.\par
% p137-b15; source PDF page 137
\blockanchor{p137-b15}
\question{3}解答如下问题:\par
% p137-b16; source PDF page 137
\blockanchor{p137-b16}
(1) 证明: $f(x, y) = \sqrt{x^2 + y^2}$ 在 $\mathbb{R}^2$ 上一致连续.\par
% p137-b17; source PDF page 137
\blockanchor{p137-b17}
(2) 证明: $h(x, y) = (x^2 + y^2)^{\frac{3}{2}}$ 在 $\mathbb{R}^2$ 上非一致连续.\par
% p137-b18; source PDF page 137
\blockanchor{p137-b18}
\question{4}设 $(x_0, y_0) \in \mathbb{R}^2$, $f(x, y)$ 在 $(x_0, y_0)$ 处连续, $g(x, y)$ 在 $(x_0, y_0)$ 处可微且 $g(x_0, y_0) = 0$, 证明: $f(x, y)g(x, y)$ 在 $(x_0, y_0)$ 处可微.\par
% p137-b19; source PDF page 137
\blockanchor{p137-b19}
\question{5}设 $f \in C^2[0,1]$, $f(0) = f(1) = 0$, $f(x)$ 在 (0,1) 上无零点, 证明: $\int_0^1 \left| \frac{f''(x)}{f(x)} \right| \, \mathrm{d}x > 4$.\par
% p137-b20; source PDF page 137
\blockanchor{p137-b20}
\question{6}设 $k > 1$, 且 $ky - \sin y + x \cos y = 0$.\par
% p137-b21; source PDF page 137
\blockanchor{p137-b21}
(1) 证明: 该方程在 $|x| < k - 1$, $|y| < +\infty$ 上确定唯一一个连续可导解 $y = y(x)$.\par
% p137-b22; source PDF page 137
\blockanchor{p137-b22}
(2) 求 $\lim_{x\to 0}\frac{y(x)}{x}$.\par
% p137-b23; source PDF page 137
\blockanchor{p137-b23}
\question{7}解答如下问题:\par
% p137-b24; source PDF page 137
\blockanchor{p137-b24}
(1) 证明: $\sum_{n=0}^{\infty} \int_{0}^{x} t^{n} \sin \pi t \, \mathrm{d}t = \int_{0}^{x} \frac{\sin \pi t}{1 - t} \, \mathrm{d}t$, 其中 $0 \leq x \leq 1$.\par
% p137-b25; source PDF page 137
\blockanchor{p137-b25}
(2) 证明: $\sum_{n=0}^{\infty} \int_{0}^{x} t^{n} \sin \pi t \, \mathrm{d}t$ 在 $[0,1]$ 上一致收敛.\par
% p138-b00; source PDF page 138
\blockanchor{p138-b00}
\question{8}设有收敛级数 $\sum_{n=1}^{\infty} a_n = S$.\par
% p138-b01; source PDF page 138
\blockanchor{p138-b01}
(1) 求  $\lim_{n\to\infty}\frac{a_{1}+2a_{2}+\cdots+na_{n}}{n}.$\par
% p138-b02; source PDF page 138
\blockanchor{p138-b02}
(2) 证明: $\sum_{n=1}^{\infty} \frac{a_1 + 2a_2 + \cdots + na_n}{n(n+1)}$ 收敛, 并求其和.\par
% p138-b03; source PDF page 138
\blockanchor{p138-b03}
\question{9}求 $I = \iiint_{D} (x + y) \, \mathrm{d}x \, \mathrm{d}y \, \mathrm{d}z$, 其中 $D$ 为 $x = 0$, $x = 1$, $x^2 + 1 = \frac{y^2}{a^2} + \frac{z^2}{b^2}$ ($a, b > 0$) 围成的区域.\par
